\scanpage{51}
\mspage{13}
\[
\mathcal H^0:\widetilde C_A\longrightarrow\widehat C_A \qquad(10)
\]
le foncteur d'injection des faisceaux dans les préfaisceaux.
Le foncteur $\mathcal H^0$ est un foncteur exact à gauche (2.3.14.II).
Ses foncteurs dérivés droitsseront notés :
\[
\mathcal H^q:\widetilde C_A\longrightarrow\widehat C_A \qquad(11)
\]
Pour tout objet $X$ de $C$ \deleted{\illegible{1}} et pour tout A-Module $F$, on a :
\[
\inserted{\mathcal H^q(F)(X)}\;\deleted{\unsure{\mathcal H^q}}\;
\deleted{H^0(\widehat C/X,\mathcal H^q(F))}=H^q(\widetilde C/X,F)\qquad(12)
\]
En effet pour $q=0$, cet isomorphisme résulte immédiatement des définitions. On l'en déduit pour tout $q$ en remarquant que le foncteur ``section\inserted{s} sur $X$'' est exact à sur la catégorie des préfaisceaux.

Soit $L:\widehat C\longrightarrow\widehat C$ le foncteur défini en (2.3.11.II).
Le foncteur $L$ commute aux limites projectives finies et par suite transforme \inserted{les}\margin{$\unsure{\rho}\quad\unsure{\rho}$} préfaisceaux d' A-modules en préfaisceaux de A-modules. Nous noterons :
\[
\check{\mathcal H}^{0}:\widehat C_A\longrightarrow\widehat C_A \qquad(13)
\]
le foncteur induit par $L$ sur les préfaisceaux de A-modules.
Le foncteur $\check{\mathcal H}^0$ est exact à \deleted{droite} gauche. Ses foncteurs dérivés droits seront notés $\check{\mathcal H}^q$. \deleted{\unsure{Les préfai}} Soit $G$ un préfaisceau de A-modules. Les préfaisceaux $\check{\mathcal H}^q(G)$ seront appelés les préfaisceaux de \uline{cohomologie de \v{C}ech} du préfaisceau $G$.

Ppour tout objet $X$ de $C$, on a un isomorphisme canonique :
\deleted{
\[
\deleted{\unsure{H^0(\widehat C, H^0(\widehat C/X,\qquad(G)))=\operatorname{Lim}}}
\]
\[
H^0(\widehat C/X,\qquad(G))=\underset R{\operatorname{Lim}}\,H^q(\widehat C/R,G)^{\unsure{C}}
\]
En effet pour $q=0$, cet isomorphisme résulte immédiatement des définitions
}
\[
\inserted{\check{\mathcal H}^q(G)(X)}\;\deleted{H^0(\widehat C/X,\check{\mathcal H}^q(G))}
=\varinjlim_R H^q(\widehat C/R,G)\qquad(14)
\]
\deleted{Ainsi qu'il} \inserted{Comme} resulte immédiatement des définitions (2.3.11.II).

(Dans cette formule la limite inductive est prise suivant la catégorie des cribles couvrants $X$).
\margin{50}
